% (c) Nikita Lisitsa, lisyarus@gmail.com, 2026

\documentclass[10pt]{beamer}

\usepackage[T2A]{fontenc}
\usepackage[russian]{babel}
\usepackage{minted}

\usepackage[most]{tcolorbox}
\tcbuselibrary{minted}

\usepackage{graphicx}
\graphicspath{ {./images/} }

\usepackage{adjustbox}

\usepackage{color}
\usepackage{soul}
\usepackage{multicol}
\usepackage{hyperref}
\usepackage{tikz}
\usetikzlibrary{arrows.meta,positioning,calc}

\usetheme{metropolis}

\definecolor{red}{rgb}{1,0,0}
\definecolor{codebg}{RGB}{29,35,49}
\setminted{fontsize=\footnotesize}

\makeatletter
\newcommand{\slideimage}[1]{
  \begin{figure}
    \begin{adjustbox}{width=\textwidth, totalheight=\textheight-2\baselineskip-2\baselineskip,keepaspectratio}
      \includegraphics{#1}
    \end{adjustbox}
  \end{figure}
}
\makeatother

\title{Язык программирования C++}
\subtitle{Лекция 16: STL, функциональные объекты, стандартные алгоритмы, iterator traits, категории итераторов, tag dispatching}
\date{2026}
\author{Никита Лисица}

\setbeamertemplate{footline}[frame number]

\begin{document}

\frame{\titlepage}

\begin{frame}[fragile]
\frametitle{Напоминание: итераторы}
\begin{itemize}
\usemintedstyle{tango}
\item Классы, позволяющиеся итерироваться по коллекции (в том числе указывать на конкретный её элемент)
\pause
\item Позволяют писать универсальный код, не зависящий от конкретного контейнера
\pause
\item По семантике и синтаксису имитируют указатели:
\usemintedstyle{lightbulb}
\begin{tcblisting}{listing engine=minted, minted language=cpp, minted options={fontsize=\scriptsize}, listing only, colback=codebg, colframe=codebg, rounded corners}
std::list<int> l;
...
auto it = l.begin();
// Сдвинуться на следующий элемент
++it;
// Прочитать текущий элемент
int x = *it;
// Записать текущий элемент
*it = 42;
\end{tcblisting}
\end{itemize}
\end{frame}

\begin{frame}[fragile]
\frametitle{Напоминание: итераторы}
\begin{itemize}
\usemintedstyle{tango}
\item Диапазон элементов задаётся парой в виде полуоткрытого интервала \mintinline{cpp}|[begin, end)|
\pause
\item \mintinline{cpp}|begin| -- итератор на первый элемент диапазона
\pause
\item \mintinline{cpp}|end| -- итератор \textit{за} последний элемент диапазона (на первый элемент \textit{после} диапазона)
\pause
\usemintedstyle{lightbulb}
\begin{tcblisting}{listing engine=minted, minted language=cpp, minted options={fontsize=\scriptsize}, listing only, colback=codebg, colframe=codebg, rounded corners}
+----+----+----+----+----+----+----+----+----+----+----+
|  4 | 47 | -5 | 19 | 100| 16 | 31 |-87 | 59 |  0 | 67 |
+----+----+----+----+----+----+----+----+----+----+----+

^                                                      ^
|                                                      |
begin                                                  end
\end{tcblisting}
\pause
\usemintedstyle{tango}
\item \mintinline{cpp}|end| часто используется для указания отсутствия элемента (напр. при поиске)
\end{itemize}
\end{frame}

\begin{frame}[fragile]
\frametitle{Функциональные объекты}
\begin{itemize}
\usemintedstyle{tango}
\item Класс, у которого перегружен оператор вызова \mintinline{cpp}|operator()|
\usemintedstyle{lightbulb}
\begin{tcblisting}{listing engine=minted, minted language=cpp, minted options={fontsize=\scriptsize}, listing only, colback=codebg, colframe=codebg, rounded corners}
struct add {
  int operator()(int x, int y) const {
    return x + y;
  }
};

add func;
int x = func(3, 4); // x = 7
\end{tcblisting}
\pause
\item \textbf{NB}: Раньше такие классы называли \textit{функторами}, но это пересекается с понятием из теории категорий, поэтому сейчас лучше говорить \textit{функциональный объект}
\end{itemize}
\end{frame}

\begin{frame}[fragile]
\frametitle{Функциональные объекты}
\begin{itemize}
\usemintedstyle{tango}
\item По смыслу похожи на функции, но мощнее\pause: позволяют хранить в себе данные
\usemintedstyle{lightbulb}
\begin{tcblisting}{listing engine=minted, minted language=cpp, minted options={fontsize=\scriptsize}, listing only, colback=codebg, colframe=codebg, rounded corners}
struct less_than {
  int reference;

  bool operator()(int x) const {
    return x < reference;
  }
};

less_than pred(10);
pred(8); // true
pred(12); // false
\end{tcblisting}
\end{itemize}
\end{frame}

\begin{frame}[fragile]
\frametitle{Предикаты}
\begin{itemize}
\usemintedstyle{tango}
\item Функциональные объекты, возвращающие \mintinline{cpp}|bool|, называют \textit{предикатами}
\pause
\item Предикат с одним аргументом -- \textit{унарный (unary predicate)}
\pause
\item Предикат с двумя аргументами -- \textit{бинарный (binary predicate)}
\pause
\usemintedstyle{lightbulb}
\begin{tcblisting}{listing engine=minted, minted language=cpp, minted options={fontsize=\scriptsize}, listing only, colback=codebg, colframe=codebg, rounded corners}
// Бинарный предикат
struct contains {
  int reference;

  bool operator()(int x, int y) const {
    return x <= reference && reference <= y;
  }
};

contains pred(10);
pred(6, 7); // false
pred(8, 12); // true
pred(12, 15); // false
\end{tcblisting}
\end{itemize}
\end{frame}

\begin{frame}[fragile]
\frametitle{Предикаты}
\begin{itemize}
\usemintedstyle{tango}
\item Функциональные объекты могут менять своё состояние при вызове
\pause
\usemintedstyle{lightbulb}
\begin{tcblisting}{listing engine=minted, minted language=cpp, minted options={fontsize=\scriptsize}, listing only, colback=codebg, colframe=codebg, rounded corners}
struct accumulator {
  int sum = 0;

  void operator()(int x) {
    sum += x;
  }
};

accumulator f;
f(10);
f(25);
f(32);

// f.sum == 67
\end{tcblisting}
\end{itemize}
\end{frame}

\begin{frame}[fragile]
\frametitle{Стандартные алгоритмы}
\begin{itemize}
\usemintedstyle{tango}
\item В STL более 100 алгоритмов, почти все -- в заголовочном файле \mintinline{cpp}|<algorithm>|
\pause
\begin{itemize}
\item Алгоритмы поиска
\pause
\item Алгоритмы, модифицирующие последовательности
\pause
\item Алгоритмы разбиения и сортировки
\pause
\item Алгоритмы на сортированных последовательностях
\pause
\item Алгоритмы работы с кучей (структурой данных)
\pause
\item Алгоритмы генерации перестановок
\pause
\item И т.д.
\end{itemize}
\pause
\item Мы рассмотрим только часть из них
\end{itemize}
\end{frame}

\begin{frame}[fragile]
\frametitle{Мини-алгоритмы}
\begin{itemize}
\usemintedstyle{tango}
\item \mintinline{cpp}|swap(T& x, T& y)| -- поменять местами два значения
\pause
\item \mintinline{cpp}|iter_swap(It p, It q)| -- поменять местами значения, на которые указывают итераторы
\pause
\item \mintinline{cpp}|min(const T& x, const T& y)| -- минимум из двух значений
\pause
\item \mintinline{cpp}|max(const T& x, const T& y)| -- максимум из двух значений
\pause
\item \mintinline{cpp}|clamp(const T& x, const T& min, const T& max)| -- ограничить значение диапазоном \mintinline{cpp}|[min, max]|
\end{itemize}
\end{frame}

\begin{frame}[fragile]
\frametitle{Min/max с компаратором}
\begin{itemize}
\usemintedstyle{tango}
\item Если стандартный алгоритм использует сравнение оператором \mintinline{cpp}|<|, у него всегда есть версия, принимающая произвольный бинарный предикат -- \textit{компаратор}
\pause
\usemintedstyle{lightbulb}
\begin{tcblisting}{listing engine=minted, minted language=cpp, minted options={fontsize=\scriptsize}, listing only, colback=codebg, colframe=codebg, rounded corners}
struct point {
  float x;
  float y;
};

struct less_by_x {
  bool operator()(const point& p1, const point& p2) const {
    return p1.x < p2.x;
  }
};

point p1{1.f, 4.f};
point p2{3.f, 2.f};

// Шаблонные параметры выводятся автоматически
std::max(p1, p2, less_by_x());
\end{tcblisting}
\end{itemize}
\end{frame}

\begin{frame}[fragile]
\frametitle{Немодифицирующие алгоритмы на последовательностях}
\begin{itemize}
\usemintedstyle{tango}
\item \mintinline{cpp}|count(It begin, It end, const T& x)| -- сколько раз в последовательности встречается значение \mintinline{cpp}|x|
\pause
\item \mintinline{cpp}|count_if(It begin, It end, UnaryPred pred)| -- сколько раз в последовательности встречаются элементы, удовлетворяющие предикату
\pause
\usemintedstyle{lightbulb}
\begin{tcblisting}{listing engine=minted, minted language=cpp, minted options={fontsize=\scriptsize}, listing only, colback=codebg, colframe=codebg, rounded corners}
struct divisible_by {
  int divisor;
  bool operator()(int x) const {
    return (x % divisor) == 0;
  }
};

std::vector<int> v;
...
std::count_if(v.begin(), v.end(), divisible_by(8));
\end{tcblisting}
\end{itemize}
\end{frame}

\begin{frame}[fragile]
\frametitle{Немодифицирующие алгоритмы}
\begin{itemize}
\usemintedstyle{tango}
\item \mintinline{cpp}|all_of(It begin, It end, UnaryPred pred)| -- правда ли, что все элементы последовательности удовлетворяют предикату
\pause
\item \mintinline{cpp}|any_of(It begin, It end, UnaryPred pred)| -- правда ли, что хотя бы один элемент последовательности удовлетворяет предикату
\pause
\item \mintinline{cpp}|none_of(It begin, It end, UnaryPred pred)| -- правда ли, что все элементы последовательности не удовлетворяют предикату
\end{itemize}
\end{frame}

\begin{frame}[fragile]
\frametitle{Алгоритмы сравнения}
\begin{itemize}
\usemintedstyle{tango}
\item \mintinline{cpp}|equal(It1 begin1, It1 end1, It2 begin2)| -- совпадают ли последовательности (должны быть одинаковой длины, поэтому не принимает \mintinline{cpp}|end2|)
\pause
\item \mintinline{cpp}|equal(It1 begin1, It1 end1, It2 begin2, It2 end2)| -- то же самое (последовательности разной длины считаются не совпадающими)
\pause
\item \mintinline{cpp}|pair<It1, It2> mismatch(It1 begin1, It1 end1, It2 begin2)| -- вернуть первое несовпадение в последовательностях
\pause
\item \mintinline{cpp}|pair<It1, It2> mismatch(It1, It1, It2, It2)| -- аналогично
\pause
\item \mintinline{cpp}|lexicographical_compare(It1, It1, It2, It2)| -- сравнить две последовательности лексикографически
\end{itemize}
\end{frame}

\begin{frame}[fragile]
\frametitle{Алгоритмы поиска}
\begin{itemize}
\usemintedstyle{tango}
\item \mintinline{cpp}|It find(It begin, It end, const T& x)| -- найти первое вхождение элемента \mintinline{cpp}|x| (или \mintinline{cpp}|end|, если такого нет)
\pause
\item \mintinline{cpp}|It find_if(It begin, It end, UnaryPred pred)| -- найти первое вхождение элемента, удовлетворяющего предикату
\pause
\item \mintinline{cpp}|It min_element(It begin, It end)| -- найти минимальный элемент последовательности
\pause
\item \mintinline{cpp}|It max_element(It begin, It end)| -- найти максимальный элемент последовательности
\pause
\item \mintinline{cpp}|pair<It, It> minmax_element(It begin, It end)| -- найти минимальный и максимальный элементы последовательности (за один проход)
\pause
\item \mintinline{cpp}|It lower/upper_bound(It begin, It end, const T& x)| -- двоичный поиск
\end{itemize}
\end{frame}

\begin{frame}[fragile]
\frametitle{Модифицирующие алгоритмы}
\begin{itemize}
\usemintedstyle{tango}
\item \mintinline{cpp}|fill(It begin, It end, const T& x)| -- заполнить последовательность копиями значения \mintinline{cpp}|x|
\pause
\item \mintinline{cpp}|generate(It begin, It end, Generator gen)| -- заполнить последовательность результатами вызова \mintinline{cpp}|gen()|
\pause
\item \mintinline{cpp}|copy(ItIn src_begin, ItIn src_end, ItOut dst_begin)| -- скопировать последовательность в последовательность, начинающуюся в \mintinline{cpp}|dst_begin|
\pause
\item \mintinline{cpp}|copy_if(ItIn, ItIn, ItOut, UnaryPred)| -- скопировать элементы, удовлетворяющие предикату
\pause
\item \mintinline{cpp}|reverse(It begin, It end)| -- обратить порядок элементов последовательности
\pause
\item \mintinline{cpp}|rotate(It begin, It middle, It end)| -- циклически сдвинуть элементы последовательности
\pause
\item \mintinline{cpp}|transform(ItIn, ItIn, ItOut, UnaryOp op)| -- как \mintinline{cpp}|copy|, но копируется результат применения \mintinline{cpp}|op(x)| к элементам
\end{itemize}
\end{frame}

\begin{frame}[fragile]
\frametitle{Алгоритмы сортировки}
\begin{itemize}
\usemintedstyle{tango}
\item \mintinline{cpp}|sort(It begin, It end)| -- отсортировать последовательность
\pause
\item \mintinline{cpp}|stable_sort(It begin, It end)| -- отсортировать стабильной сортировкой (равные элементы сохраняют порядок)
\pause
\item \mintinline{cpp}|partial_sort(It begin, It middle, It end)| -- сдвинуть младшие \mintinline{cpp}|middle - begin| элементов последовательности в начало
\pause
\item Есть версии, принимающие компаратор
\end{itemize}
\end{frame}

\begin{frame}[fragile]
\frametitle{Пример реализации алгоритма}
\begin{itemize}
\usemintedstyle{lightbulb}
\begin{tcblisting}{listing engine=minted, minted language=cpp, minted options={fontsize=\scriptsize}, listing only, colback=codebg, colframe=codebg, rounded corners}
template <typename InputIt, typename OutputIt>
OutputIt copy(InputIt src_begin, InputIt src_end,
  OutputIt dst_begin)
{
  for (; src_begin != src_end;)
    *dst_begin++ = *src_begin++;

  return dst_begin;
}
\end{tcblisting}
\pause
\usemintedstyle{tango}
\item \mintinline{cpp}|*p++ = *q++;| -- классический способ скопировать последовательность
\end{itemize}
\end{frame}

\begin{frame}[fragile]
\frametitle{Пример реализации алгоритма}
\begin{itemize}
\usemintedstyle{lightbulb}
\begin{tcblisting}{listing engine=minted, minted language=cpp, minted options={fontsize=\scriptsize}, listing only, colback=codebg, colframe=codebg, rounded corners}
template <typename InputIt, typename OutputIt,
  typename UnaryPred>
OutputIt copy(InputIt src_begin, InputIt src_end,
  OutputIt dst_begin, UnaryPred pred)
{
  for (; src_begin != src_end; ++src_begin)
    if (pred(*src_begin))
      *dst_begin++ = *src_begin;

  return dst_begin;
}
\end{tcblisting}
\pause
\usemintedstyle{tango}
\item На \href{https://en.cppreference.com}{\nolinkurl{cppreference.com}} почти к любому алгоритму есть примеры реализации
\end{itemize}
\end{frame}

\begin{frame}[fragile]
\frametitle{Iterator traits}
\begin{itemize}
\item Хотим реализовать быструю сортировку
\usemintedstyle{lightbulb}
\begin{tcblisting}{listing engine=minted, minted language=cpp, minted options={fontsize=\scriptsize}, listing only, colback=codebg, colframe=codebg, rounded corners}
template <typename It>
void quick_sort(It begin, It end)
{
  // Выбираем опорную точку для сортировки
  const T& pivot = ...;
}
\end{tcblisting}
\pause
\usemintedstyle{tango}
\item Откуда нам взять тип \mintinline{cpp}|T|?
\end{itemize}
\end{frame}

\begin{frame}[fragile]
\frametitle{Iterator traits}
\begin{itemize}
\item Добавим тип элемента в сам класс итератора:
\usemintedstyle{lightbulb}
\begin{tcblisting}{listing engine=minted, minted language=cpp, minted options={fontsize=\scriptsize}, listing only, colback=codebg, colframe=codebg, rounded corners}
template <typename T>
class vector {
  class iterator {
    using value_type = T;
  };
};

template <typename It>
void quick_sort(It begin, It end)
{
  // Выбираем опорную точку для сортировки
  const It::value_type& pivot = ...;
}
\end{tcblisting}
\pause
\usemintedstyle{tango}
\item Ошибка компиляции: компилятор не понимает, \mintinline{cpp}|It::value_type| это тип или статическая переменная
\end{itemize}
\end{frame}

\begin{frame}[fragile]
\frametitle{Iterator traits}
\begin{itemize}
\usemintedstyle{tango}
\item Добавим волшебное слово \mintinline{cpp}|typename|:
\usemintedstyle{lightbulb}
\begin{tcblisting}{listing engine=minted, minted language=cpp, minted options={fontsize=\scriptsize}, listing only, colback=codebg, colframe=codebg, rounded corners}
template <typename T>
class vector {
  class iterator {
    using value_type = T;
  };
};

template <typename It>
void quick_sort(It begin, It end)
{
  // Выбираем опорную точку для сортировки
  const typename It::value_type& pivot = ...;
}
\end{tcblisting}
\pause
\usemintedstyle{tango}
\item Теперь всё компилируется, но не работает если \mintinline{cpp}|It| это указатель :(
\end{itemize}
\end{frame}

\begin{frame}[fragile]
\frametitle{Iterator traits}
\begin{itemize}
\usemintedstyle{tango}
\item Решение: \textit{iterator traits (от англ. trait -- свойство, признак, черта)}
\pause
\item Стандартный шаблонный класс, параметризующийся типом итератора, и хранящий некую информацию о нём
\pause
\item По умолчанию достаёт все данные из самого типа итератора:
\usemintedstyle{lightbulb}
\begin{tcblisting}{listing engine=minted, minted language=cpp, minted options={fontsize=\scriptsize}, listing only, colback=codebg, colframe=codebg, rounded corners}
template <typename Iterator>
struct iterator_traits {
  using value_type = typename Iterator::value_type;
  using reference_type = typename Iterator::reference_type;
  ...
};
\end{tcblisting}
\end{itemize}
\end{frame}

\begin{frame}[fragile]
\frametitle{Iterator traits}
\begin{itemize}
\usemintedstyle{tango}
\item Есть частичная специализация для указателей:
\usemintedstyle{lightbulb}
\begin{tcblisting}{listing engine=minted, minted language=cpp, minted options={fontsize=\scriptsize}, listing only, colback=codebg, colframe=codebg, rounded corners}
template <typename T>
struct iterator_traits<T*> {
  using value_type = T;
  using reference_type = T&;
  ...
};

template <typename T>
struct iterator_traits<const T*> {
  using value_type = T const;
  using reference_type = T const&;
  ...
};
\end{tcblisting}
\end{itemize}
\end{frame}

\begin{frame}[fragile]
\frametitle{Iterator traits}
\begin{itemize}
\usemintedstyle{tango}
\item Решение нашей проблемы:
\usemintedstyle{lightbulb}
\begin{tcblisting}{listing engine=minted, minted language=cpp, minted options={fontsize=\scriptsize}, listing only, colback=codebg, colframe=codebg, rounded corners}
template <typename It>
void quick_sort(It begin, It end)
{
  // Работает как для указателей, так и для vector::iterator
  using traits = std::iterator_traits<It>;
  const typename traits::value_type& pivot = ...;
}
\end{tcblisting}
\pause
\usemintedstyle{tango}
\item \textbf{NB}: во многих случаях вместо \mintinline{cpp}|std::iterator_traits| проще использовать вывод типов (\mintinline{cpp}|auto|):
\usemintedstyle{lightbulb}
\begin{tcblisting}{listing engine=minted, minted language=cpp, minted options={fontsize=\scriptsize}, listing only, colback=codebg, colframe=codebg, rounded corners}
template <typename It>
void quick_sort(It begin, It end)
{
  const auto& pivot = *begin;
}
\end{tcblisting}
\end{itemize}
\end{frame}

\begin{frame}[fragile]
\frametitle{Алгоритмы на итераторах}
\begin{itemize}
\usemintedstyle{tango}
\item \mintinline{cpp}|next(it)| -- вернуть следующий итератор, не меняя сам \mintinline{cpp}|it|
\pause
\item \mintinline{cpp}|next(it, n)| -- вернуть копию, сдвинутую на \mintinline{cpp}|n| элементов
\pause
\item \mintinline{cpp}|advance(it, n)| -- сдвинуть \mintinline{cpp}|it| (принимается по ссылке) на \mintinline{cpp}|n| элементов
\pause
\item \mintinline{cpp}|distance(it1, it2)| -- на сколько нужно сдвинуть \mintinline{cpp}|it1|, чтобы получить \mintinline{cpp}|it2|
\pause
\item Как быстро работают эти алгоритмы? \pause Зависит от конкретного контейнера: \(O(1)\) для вектора, \(O(N)\) для списка
\pause
\item Как написать разную реализацию алгоритмов для разных контейнеров? \pause \(\Longrightarrow\) \textit{категории итераторов}
\end{itemize}
\end{frame}

\begin{frame}[fragile]
\frametitle{Категории итераторов}
\begin{itemize}
\usemintedstyle{tango}
\item Категория итератора определяет, какие операции этот итератор умеет выполнять (и эффективно)
\pause
\begin{itemize}
\item \textit{Forward iterator}: только \mintinline{cpp}|*it| и \mintinline{cpp}|++it|
\pause
\item \textit{Bidirectional iterator}: ещё \mintinline{cpp}|--it|
\pause
\item \textit{Random access iterator}: \mintinline{cpp}|it += n| и \mintinline{cpp}|it -= n|
\end{itemize}
\pause
\item Итераторы по вектору или деке -- random access
\pause
\item Итераторы по списку или set/map -- bidirectional
\pause
\item Итераторы по односвязному списку -- forward
\end{itemize}
\end{frame}

\begin{frame}[fragile]
\frametitle{Категории итераторов}
\begin{itemize}
\usemintedstyle{tango}
\item В стандартной библитеке есть специальные классы-теги для определения категорий итераторов:
\usemintedstyle{lightbulb}
\begin{tcblisting}{listing engine=minted, minted language=cpp, minted options={fontsize=\scriptsize}, listing only, colback=codebg, colframe=codebg, rounded corners}
struct forward_iterator_tag {};
struct bidirectional_iterator_tag : forward_iterator_tag {};
struct random_access_iterator_tag : bidirectional_iterator_tag{};
\end{tcblisting}
\pause
\usemintedstyle{tango}
\item Они наследуются друг от друга, создавая иерархию категорий
\end{itemize}
\end{frame}

\begin{frame}[fragile]
\frametitle{Категории итераторов}
\begin{itemize}
\usemintedstyle{tango}
\item В \mintinline{cpp}|iterator_traits| (и в самих классах итераторов) указывается их категория:
\usemintedstyle{lightbulb}
\begin{tcblisting}{listing engine=minted, minted language=cpp, minted options={fontsize=\scriptsize}, listing only, colback=codebg, colframe=codebg, rounded corners}
template <typename T>
class list {
  class iterator {
    using iterator_category = bidirectional_iterator_tag;
  };
};

template <typename T>
struct iterator_traits<T*> {
  using iterator_category = random_access_iterator_tag;
};
\end{tcblisting}
\end{itemize}
\end{frame}

\begin{frame}[fragile]
\frametitle{Tag dispatching}
\begin{itemize}
\usemintedstyle{tango}
\item Как реализовать \mintinline{cpp}|std::next(it, n)| эффективно для разных итераторов? \pause Использовать перегрузку функций:
\usemintedstyle{lightbulb}
\begin{tcblisting}{listing engine=minted, minted language=cpp, minted options={fontsize=\scriptsize}, listing only, colback=codebg, colframe=codebg, rounded corners}
template <typename It>
It next(It it, int n, std::bidirectional_iterator_tag) {
  for (; n > 0; --n)
    ++it;
  return it;
}

template <typename It>
It next(It it, int n, std::random_access_iterator_tag) {
  it += n;
  return it;
}
\end{tcblisting}
\end{itemize}
\end{frame}

\begin{frame}[fragile]
\frametitle{Tag dispatching}
\begin{itemize}
\usemintedstyle{tango}
\usemintedstyle{lightbulb}
\begin{tcblisting}{listing engine=minted, minted language=cpp, minted options={fontsize=\scriptsize}, listing only, colback=codebg, colframe=codebg, rounded corners}
template <typename It>
It next(It it, int n) {
  using traits = std::iterator_traits<It>;
  return next(it, n, typename traits::iterator_category{});
}
\end{tcblisting}
\pause
\item Такой подход -- перегрузка функции на основе дополнительного параметра-тега -- называется \textit{tag dispatching}, и очень распространён в шаблонном коде
\end{itemize}
\end{frame}

\begin{frame}[fragile]
\frametitle{Оператор \texttt{-->}}
\begin{itemize}
\usemintedstyle{tango}
\item Цикл вида \mintinline{cpp}|for (; n > 0; --n)| можно записать смешнее:
\usemintedstyle{lightbulb}
\begin{tcblisting}{listing engine=minted, minted language=cpp, minted options={fontsize=\scriptsize}, listing only, colback=codebg, colframe=codebg, rounded corners}
while (n --> 0) {
  ...
}
\end{tcblisting}
\pause
\usemintedstyle{tango}
\item Это не отдельный оператор, а комбинация постфиксного декремента \mintinline{cpp}|n--| и сравнения \mintinline{cpp}|(n--) > 0|
\pause
\item Такой цикл выполнит ровно \mintinline{cpp}|n| итераций :)
\end{itemize}
\end{frame}

\end{document}